\documentclass[11pt,a4paper]{article}

% --- Basics ---
\usepackage[utf8]{inputenc}
\usepackage[T1]{fontenc}
\usepackage[ngerman]{babel}
\usepackage{lmodern}
\usepackage{microtype}
\usepackage{geometry}
\geometry{margin=2.6cm}

% --- Math ---
\usepackage{amsmath,amssymb,amsthm,mathtools}
\usepackage{bm}

% --- Links ---
\usepackage{hyperref}

% --- Theorem environments (deutsch) ---
\theoremstyle{plain}
\newtheorem{satz}{Satz}
\newtheorem{lemma}{Lemma}
\newtheorem{korollar}{Korollar}

\theoremstyle{definition}
\newtheorem{definition}{Definition}
\newtheorem{axiom}{Axiom}
\newtheorem{conjecture}{Vermutung}

\theoremstyle{remark}
\newtheorem{bemerkung}{Bemerkung}

% --- Convenience ---
\newcommand{\C}{\mathbb{C}}
\newcommand{\R}{\mathbb{R}}
\newcommand{\Hh}{\mathbb{H}}
\newcommand{\Okt}{\mathbb{O}}
\newcommand{\ii}{\mathrm{i}}
\newcommand{\dd}{\mathrm{d}}
\newcommand{\norm}[1]{\left\lVert #1 \right\rVert}

\title{%
Die kritische Achse als Dualit\"ats-Fixpunkt\\
in den normierten Divisionsalgebren $\R,\C,\Hh,\Okt$\\
(Paper A: mathematischer Satz + klar markierte Vermutung zur RH)
}
\author{Thomas Hoffbauer\\Bamberg, Deutschland}
\date{12. Februar 2026}

\begin{document}
\maketitle

\begin{abstract}
Wir isolieren einen rein mathematischen Kern: In jeder normierten Divisionsalgebra
$\mathbb{A}\in\{\R,\C,\Hh,\Okt\}$ erzwingt die Norm-Dualit\"atsbedingung
$\|x\|=\|1-x\|$ f\"ur ein Element $x$ in einem komplexen Slice $\C_u\subset\mathbb{A}$
notwendig die Fixpunktlage $\Re(x)=1/2$.
Dieser Satz ist elementar, aber strukturell bedeutsam: Er identifiziert die ``Einhalb'' als
Symmetrieachse einer unitarit\"atsvertr\"aglichen Dualit\"at, unabh\"angig von Physik.

Die Verbindung zur Riemannschen Vermutung wird strikt getrennt: Wir formulieren eine
Vermutung, dass die nichttrivialen Nullstellen der (vollendeten) Zetafunktion genau solche
universell dualit\"atsstabilen Resonanzen unter nat\"urlichen Lifts in $\Hh$ bzw.\ $\Okt$ sind.
Damit wird ``Satz + Vermutung'' zu einem klaren, logisch transparenten Programm.
\end{abstract}

\section{Motivation und Claim}
Dieses Paper macht \emph{keine} physikalische Behauptung \"uber Strahlungsspektren.
Der Claim ist rein algebraisch-geometrisch:

\begin{quote}
\emph{Wenn ein (liftbarer) komplexer Parameter in einer normierten Divisionsalgebra die Norm-Dualit\"at
$x\mapsto 1-x$ erf\"ullt, dann liegt sein skalarer Anteil zwingend auf $1/2$.}
\end{quote}

Das ist die mathematische Form dessen, was im fr\"uheren Manuskript als ``Hurwitz-Imperativ''
auftauchte, nun jedoch als sauberer Satz mit expliziten Voraussetzungen.



\subsection*{Historische Notiz: von Hamilton zu den Oktonionen}
Die vier normierten Divisionsalgebren $\mathbb{R},\mathbb{C},\mathbb{H},\mathbb{O}$ sind nicht nur ein
klassifikatorisches Resultat (Hurwitz), sondern auch ein historischer Pfad:
\begin{itemize}
\item \textbf{1843:} \emph{William Rowan Hamilton} entdeckt die Quaternionen $\mathbb{H}$ (ber\"uhmt ist die Notiz
      am Brougham/Broom Bridge in Dublin).\footnote{\url{https://www.maths.tcd.ie/pub/HistMath/People/Hamilton/Letters/BroomeBridge.html}}
\item \textbf{1843/44:} \emph{John T.\ Graves} (ein Studienfreund Hamiltons) findet kurz danach die Oktonionen
      $\mathbb{O}$; \emph{Arthur Cayley} publiziert sp\"ater eine systematische Darstellung (daher auch ``Cayley-Zahlen'').\footnote{\url{https://mathshistory.st-andrews.ac.uk/Biographies/Graves_John/}}\footnote{\url{https://www.ams.org/bull/2002-39-02/S0273-0979-01-00934-X/}}
\end{itemize}

Diese Linie ist f\"ur uns konzeptionell wichtig: In $\mathbb{H}$ und $\mathbb{O}$ existieren viele
\emph{komplexe Slices} (``eingebettete $\mathbb{C}$-Ebenen''), und damit die M\"oglichkeit, analytische
Strukturen aus $\mathbb{C}$ in h\"ohere Divisionsalgebren zu \emph{liften}. Genau dieses ``Lift-Programm''
ist die Stelle, an der eine nichttriviale, symmetriegetriebene Selektionshypothese f\"ur Nullstellen ansetzen kann.

\section{Normierte Divisionsalgebren und komplexe Slices}

\begin{definition}[Normierte Divisionsalgebra]
Eine (reelle) Algebra $\mathbb{A}$ hei\ss t \emph{normierte Divisionsalgebra}, wenn es eine Norm
$\|\cdot\|$ gibt mit \emph{Kompositionsgesetz} $\|xy\|=\|x\|\,\|y\|$ f\"ur alle $x,y\in\mathbb{A}$
und jedes $x\neq 0$ invertierbar ist.
Nach Hurwitz gilt $\mathbb{A}\in\{\R,\C,\Hh,\Okt\}$.
\end{definition}

\begin{definition}[Komplexer Slice]
Sei $\mathbb{A}\in\{\Hh,\Okt\}$ und $u\in\mathrm{Im}(\mathbb{A})$ mit $u^2=-1$.
Dann ist
\[
\C_u := \{\,\sigma + \gamma u : \sigma,\gamma\in\R\,\}\subset\mathbb{A}
\]
ein Unterraum, der als Algebra isomorph zu $\C$ ist (``complex slice'').
\end{definition}

\begin{definition}[Slice-Einbettung]
Die kanonische Einbettung in den Slice sei
\[
\iota_u:\C\to\mathbb{A},\qquad \iota_u(\sigma+\ii\gamma)=\sigma+\gamma u.
\]
\end{definition}

\begin{bemerkung}
Die Norm in $\mathbb{A}$ reduziert sich auf dem Slice auf den euklidischen Betrag:
\[
\|\iota_u(\sigma+\ii\gamma)\|^2 = \sigma^2+\gamma^2 = |\sigma+\ii\gamma|^2.
\]
Dies ist elementar; die inhaltliche Frage liegt nicht in dieser Identit\"at, sondern darin,
\emph{wann} und \emph{f\"ur welche} Parameter eine Dualit\"ats-Normbedingung physikalisch oder
zahlentheoretisch motiviert ist.
\end{bemerkung}

\section{Dualit\"ats-Normfunktional und Fixpunkt}

\begin{definition}[Dualit\"ats-Funktional]
F\"ur $x\in\mathbb{A}$ definieren wir
\[
\Phi(x) := \big|\|x\|^2-\|1-x\|^2\big|.
\]
Auf einem Slice schreiben wir auch $\Phi_u(s):=\Phi(\iota_u(s))$.
\end{definition}

\begin{lemma}[Fixpunkt der Norm-Dualit\"at]\label{lem:fixpoint}
Sei $\mathbb{A}\in\{\R,\C,\Hh,\Okt\}$ und $x=\sigma+\gamma u\in\C_u\subset\mathbb{A}$.
Wenn $\Phi(x)=0$, dann gilt $\sigma=\tfrac12$.
\end{lemma}

\begin{proof}
Auf dem Slice gilt $\|x\|^2=\sigma^2+\gamma^2$ und $\|1-x\|^2=(1-\sigma)^2+\gamma^2$.
Aus $\Phi(x)=0$ folgt $\sigma^2+\gamma^2=(1-\sigma)^2+\gamma^2$, also $2\sigma=1$.
\end{proof}

\begin{satz}[Kritische Achse als Dualit\"ats-Fixpunkt]\label{thm:critical_axis}
In jeder normierten Divisionsalgebra $\mathbb{A}\in\{\R,\C,\Hh,\Okt\}$ liegt jedes Element
$x\in\C_u$ mit $\|x\|=\|1-x\|$ auf der Fixpunktachse $\Re(x)=1/2$.
\end{satz}

\begin{proof}
Unmittelbar aus Lemma~\ref{lem:fixpoint}.
\end{proof}

\begin{bemerkung}[Positionierung und Neuheit]
Satz~\ref{thm:critical_axis} ist elementar: Er identifiziert die Fixpunktachse der involutiven Dualität
$x\mapsto 1-x$ unter einer Normsymmetrie.
Die \emph{Neuigkeitsbehauptung} dieses Manuskripts liegt daher nicht in der algebraischen Rechnung,
sondern in (i) der Definition eines \emph{universellen} Stabilitätsbegriffs über Quaternionen-/Oktonionen-Slices
(Definition~\ref{def:universal}) und (ii) der klar als \emph{Vermutung} markierten These,
dass die nichttrivialen Zeta-Nullstellen genau durch eine solche Universality- bzw.\ Lift-Stabilität ausgezeichnet sind
(Vermutung~\ref{conj:rh}).
\end{bemerkung}



\section{Universelle Stabilit\"at \"uber Slices}
Das Theorem ist elementar, aber sein Nutzen entsteht, wenn man eine \emph{Universality}-Annahme
formuliert: dass bestimmte ``Resonanzen'' gegen Wahl des Slices stabil sind.

\begin{definition}[Universelle Dualit\"atsstabilit\"at]\label{def:universal}
Sei $\mathbb{A}\in\{\Hh,\Okt\}$ und $\mathcal{U}$ eine Klasse zul\"assiger imagin\"arer Einheiten
$u$ (z.B.\ alle $u$ mit $u^2=-1$ oder eine durch Strukturbedingungen eingeschr\"ankte Teilmenge).
Ein Parameter $s\in\C$ hei\ss t \emph{universell dualit\"atsstabil} (bez\"uglich $\mathcal{U}$), wenn
\[
\Phi(\iota_u(s))=0\quad\text{f\"ur alle }u\in\mathcal{U}.
\]
\end{definition}

\begin{korollar}[Universality $\Rightarrow$ kritische Linie]
Jeder universell dualit\"atsstabile Parameter $s$ erf\"ullt $\Re(s)=1/2$.
\end{korollar}

\subsection*{Wahl der Slice-Klasse $\mathcal U$ (kanonisch über Symmetrie)}
Die Definition lässt $\mathcal U$ bewusst offen, weil hier die inhaltliche Last liegt:
\emph{Welche} Slices sind ``physikalisch/mathematisch zulässig''?
Ein kanonischer Kandidat ist die volle Menge der imaginären Einheiten
\[
\mathcal U_{\mathrm{can}}=\{u\in\mathrm{Im}(\mathbb A): \|u\|=1,\ u^2=-1\}.
\]
Diese Wahl ist durch Symmetrie motiviert:
In $\Hh$ wirken innere Automorphismen (Isometrien der Kompositionsnorm) transitiv auf den Einheitenkreis
im imaginären Raum (Rotationen der Imaginärachse), und in $\Okt$ ist die Automorphismengruppe
$\mathrm{Aut}(\Okt)\cong G_2$ ein natürlicher Symmetrieanbieter, der die Norm erhält.
Damit ist $\mathcal U_{\mathrm{can}}$ nicht ``maßgeschneidert'', sondern der naheliegende Orbit eines
Norm-erhaltenden Symmetrieprinzips.

\subsection*{Nontriviality: Was ist hier mehr als $\Re(s)=1/2$?}
Auf den ersten Blick scheint Definition~\ref{def:universal} schwach, weil das Norm-Funktional $\Phi(\iota_u(s))$
für die minimale Slice-Einbettung nur von $(\Re(s),\Im(s))$ und nicht von der Richtung $u$ abhängt.
Der \emph{nichttriviale} Gehalt beginnt dort, wo man die Universality nicht als bloßes
``für alle $u$'' der gleichen Gleichung versteht, sondern als Stabilitätsprinzip innerhalb einer
\emph{Lift-Theorie}:
\begin{itemize}
\item In normierten Divisionsalgebren existiert eine Kompositionsnorm, die solche Dualitäts-/Fixpunktformulierungen
      überhaupt global sinnvoll macht; in nichtkompositionellen Algebren bricht diese Struktur.
\item Die programmatische Vermutung (Abschnitt~\ref{sec:conj}) behauptet nicht nur $\Re(\rho)=1/2$,
      sondern dass die Zeta-Nullstellen als \emph{universelle Resonanzen} durch ein natürliches Lift-/Stabilitätsprinzip
      ausgezeichnet sind. Das ist eine starke Selektionshypothese, nicht eine Umbenennung der kritischen Linie.
\item Kontrast: Für eine generische holomorphe Funktion $f$ ohne Dualitätssymmetrie $f(s)\leftrightarrow f(1-s)$
      ist eine solche universelle Dualitätsstabilität typischerweise \emph{leer} bzw.\ nicht mit der Nullstellenstruktur
      kompatibel; die Vermutung ist also falsifizierbar.
\end{itemize}

\begin{lemma}[Kontrastlemma: Dualitätsstabilität ist nicht generisch]\label{lem:contrast}
Sei $f:\C\to\C$ holomorph mit diskreter Nullstellenmenge $Z(f)$.
Die Bedingung
\[
|\,\rho\,|=|\,1-\rho\,|\qquad(\text{äquivalent: }\Phi(\iota_u(\rho))=0)
\]
erzwingt für jede Nullstelle $\rho\in Z(f)$ die Fixpunktlage $\Re(\rho)=\tfrac12$.
Ohne eine zusätzliche Dualitäts- bzw.\ Stabilitätsstruktur ist diese Lagebedingung jedoch
\emph{nicht} typischerweise erfüllt: Schon in einfachen Familien hängt sie von einem
echten Parameterconstraint ab.
\end{lemma}

\begin{proof}
Die Implikation $|\,\rho\,|=|\,1-\rho\,|\Rightarrow \Re(\rho)=\tfrac12$ ist die Fixpunktrechnung aus
Lemma~\ref{lem:fixpoint}. Für die Nicht-Generizität genügt ein elementares Beispiel:
Für $f_a(s)=s-a$ ist $Z(f_a)=\{a\}$, und die Dualitätsstabilität $|a|=|1-a|$ gilt genau dann,
wenn $\Re(a)=\tfrac12$. Für \emph{generische} Parameter $a$ ist sie also verletzt.
\end{proof}

\begin{bemerkung}[Beispiele und Intuition]\label{rem:examples}
\begin{enumerate}
\item \textbf{Lineares Beispiel.} $f_a(s)=s-a$ besitzt eine Nullstelle bei $a$; Dualitätsstabilität
      selektiert $\Re(a)=1/2$ als nichttriviale Einschränkung (Lemma~\ref{lem:contrast}).
\item \textbf{Vakuum-Beispiel.} $f(s)=e^s$ hat keine Nullstellen; jede ``Nullstellenbehauptung'' ist hier
      leer. Das unterstreicht: Erst die Kombination aus (i) Nullstellenstruktur und (ii) Stabilitätsaxiom
      erzeugt Inhalt.
\item \textbf{Dualität ohne Achs-Zwang (Mini-Beispiel).} 
      Allein aus der Spiegel-Symmetrie der Nullstellenmenge folgt \emph{kein} Achs-Zwang.
      Setze f\"ur ein beliebiges $a\in\C$ die Funktion
      \[
      f_a(s):=(s-a)\,(s-(1-a)).
      \]
      Dann ist $Z(f_a)=\{a,\,1-a\}$ zwar invariant unter $s\mapsto 1-s$, aber f\"ur generisches $a$ gilt
      $\Re(a)\neq\tfrac12$ und damit liegen die Nullstellen \emph{nicht} auf der Fixpunktachse.
      Die Fixpunktlage entsteht erst, wenn zus\"atzlich eine Norm-/Unitaritätsforderung der Form
      $\Phi(\iota_u(\rho))=0$ \emph{an den Nullstellen selbst} gefordert wird (vgl.\ Lemma~\ref{lem:fixpoint}). 

\end{enumerate}
\end{bemerkung}

\subsection*{Arbeitsdefinition eines Lift-Kandidaten (Slice-Regular / Stem-Framework)}
Um die oben genannte Nichttrivialität formal angreifbar zu machen, verwenden wir als \emph{Working Hypothesis}
eine Standardkonstruktion aus der Theorie der slice-regularen Funktionen.
Sei $F:\C\to\C$ holomorph und \emph{intrinsisch} (``stem''), d.\,h.
\[
F(\overline{z})=\overline{F(z)}.
\]
Schreibe $F(z)=\alpha(x,y)+\ii\,\beta(x,y)$ für $z=x+\ii y$.
Dann definieren wir für $\mathbb A\in\{\Hh,\Okt\}$ die zugehörige Slice-Funktion
\[
\mathrm{Lift}_{\mathbb A}[F](x+y\,u)\;:=\;\alpha(x,y)+u\,\beta(x,y),\qquad u\in\mathcal U.
\]
Diese Konstruktion ist kanonisch, weil sie (i) auf jedem komplexen Slice mit $F$ kompatibel ist,
(ii) die Kompositionsnorm respektiert und (iii) die Wahl von $\mathcal U$ über Symmetrie erlaubt.
In v1.1 ist dies eine \emph{Arbeitsdefinition}; eine strengere Fassung (funktionalanalytische Domäne,
Regularität, eindeutige Fortsetzung) gehört zu den offenen Punkten.

\subsection*{Computational Test Plan (Operationalisierung der Vermutung)}
Die Vermutung wird nur dann wissenschaftlich fruchtbar, wenn sie operationalisierbar ist.
Ein minimaler numerischer Plan (ohne Physik) lautet:
\begin{enumerate}
\item Fixiere eine endliche, symmetrische Stichprobe $\{u_k\}_{k=1}^K\subset\mathcal U_{\mathrm{can}}$
      (z.\,B.\ durch gleichmäßig verteilte Punkte auf $S^2$ bzw.\ $S^6$ oder ein sphärisches $t$-Design).
\item Implementiere den Lift $\mathrm{Lift}_{\mathbb A}$ für Kandidaten $F$ (z.\,B.\ $\xi$ oder geeignete Modifikationen)
      und evaluiere $\mathrm{Lift}_{\mathbb A}[F](x+y\,u_k)$.
\item Prüfe für bekannte numerische Nullstellenkandidaten $\rho_n=\tfrac12+\ii\gamma_n$:
      (a) Robustheit der Nullstellenbedingung unter Variation von $u_k$ (bis auf numerische Toleranz),
      (b) Konsistenz mit einer Dualitäts-/Stabilitätsbedingung, die $\Phi=0$ erzwingt.
\item Ein klares Gegenbeispiel wäre: systematische, $u$-abhängige Drift der Nullstellenbedingung, die nicht durch
      Numerik erklärbar ist, oder Stabilität nur für speziell ausgewählte $u$.
\end{enumerate}
Dieser Plan ersetzt keinen Beweis, macht aber die ``Universality'' als Programm testbar.


\section{Vermutung zur Riemannschen Vermutung (klar markiert)}\label{sec:conj}
Wir trennen streng zwischen dem Satzteil (bewiesen) und der Br\"ucke zur Zetafunktion (offen).

\begin{conjecture}[Division-Algebra Universality der nichttrivialen Nullstellen]\label{conj:rh}
Seien $\rho$ nichttriviale Nullstellen der vollendeten Zetafunktion (oder $\xi$-Funktion).
Dann sind diese $\rho$ universell dualit\"atsstabil im Sinn von
Definition~\ref{def:universal} f\"ur eine nat\"urliche Klasse $\mathcal{U}$ von Slices, die durch den
Lift/Analytische Fortsetzung der Riemann-Funktion in $\Hh$ bzw.\ $\Okt$ kanonisch bestimmt ist.
\end{conjecture}

\begin{korollar}[Programm: Satz + Vermutung $\Rightarrow$ RH]
Gilt die obige Vermutung, dann folgt f\"ur jede nichttriviale Nullstelle $\rho$:
$\Re(\rho)=1/2$.
\end{korollar}

\section{Offene Punkte (was zu zeigen w\"are)}
\begin{itemize}
\item Eine kanonische (nichttriviale) Definition des ``Lifts'' der analytisch fortgesetzten Zeta-/
$\xi$-Funktion in $\Hh$ und $\Okt$ (z.B.\ Slice-regular oder Stem-function Framework).
\item Eine begr\"undete Wahl der Slice-Klasse $\mathcal{U}$, die die Zahlentheorie (Funktionalgleichung)
und die Algebra (Automorphismen/Isometrien der Norm) zugleich reflektiert.
\item Der Nachweis, dass die nichttrivialen Nullstellen die Universality-Bedingung erf\"ullen.
\end{itemize}

\section*{Hinweis zum Folgepaper (Paper B)}
Der physikalische Teil (Strahlungsmodulation, Template-Matching, Metrologie) wird bewusst in einem
separaten Manuskript behandelt, um die logische Struktur dieses Papers rein mathematisch zu halten.

\end{document}
